\documentclass{article}
\usepackage[utf8]{inputenc}
\usepackage[T2A]{fontenc}
\usepackage{helvet}
\usepackage{geometry}
\usepackage{xcolor}
\usepackage{eso-pic}
\usepackage{fancyhdr}
\usepackage{parskip}
\usepackage{microtype}
\usepackage[english, ukrainian]{babel}
\usepackage{url}
\usepackage{amsmath}
\usepackage{amssymb}
\usepackage{amsthm}
\usepackage{longtable}
\usepackage{booktabs}
\usepackage{array}
\usepackage{hyperref}
\usepackage{titlesec}

% Define brand colors
\definecolor{primary}{HTML}{293757}
\definecolor{footergray}{gray}{0.65}

\geometry{a4paper,left=3cm,right=3cm,top=3cm,bottom=3cm}
\renewcommand{\familydefault}{\sfdefault}
\setcounter{secnumdepth}{3}

% Font shapes for T2A/Helvetica (mapping to cmss for Cyrillic support)
\DeclareFontFamily{T2A}{phv}{}
\DeclareFontShape{T2A}{phv}{m}{n}{<->ssub * cmss/m/n}{}
\DeclareFontShape{T2A}{phv}{bx}{n}{<->ssub * cmss/bx/n}{}
\DeclareFontShape{T2A}{phv}{m}{it}{<->ssub * cmss/m/it}{}
\DeclareFontShape{T2A}{phv}{bx}{it}{<->ssub * cmss/bx/it}{}

\titleformat{\section}
  {\color{primary}\Huge\bfseries}{\thesection.}{1em}{}
\titlespacing*{\section}{0pt}{30pt}{12pt}
\titleformat{\subsection}
  {\color{primary}\LARGE\bfseries}{\thesubsection.}{1em}{}
\titlespacing*{\subsection}{0pt}{20pt}{8pt}
\titleformat{\subsubsection}
  {\color{primary}\Large\bfseries}{\thesubsubsection.}{1em}{}
\titlespacing*{\subsubsection}{0pt}{10pt}{6pt}

\fancyhf{}
\fancyhead[R]{\small\color{footergray} 29 червня 2026}
\fancyfoot[L]{\small\color{footergray} Техноробочий проект ERP/1. Модуль Істотність}
\fancyfoot[R]{\small\color{footergray}\thepage}

\newcommand\HUGE[1]{\fontsize{38}{38}\selectfont #1}
\renewcommand{\headrulewidth}{0pt}
\renewcommand{\footrulewidth}{0.4pt}
\renewcommand{\footrule}{\color{footergray}\hrule width \textwidth height 0.4pt}

\begin{document}
\pagestyle{fancy}

\begin{titlepage}
\AddToShipoutPictureBG*{\AtPageLowerLeft{\color{primary}\rule{\paperwidth}{\paperheight}}}
\color{white}\thispagestyle{empty}\vspace*{3cm}\centering
  {\Huge \textbf{ERP/1: Істотність} \par} \vspace{1.5cm}
  {\HUGE \textbf{Техноробочий проект} \par} \vspace{1.5cm}
  {\LARGE \textbf{на створення та впровадження \\ локальної LLM-інфраструктури \\ для аналізу та оцінки істотності інформації} \par} \vspace{1.5cm}
  {\large Згідно з Постановою КМУ № 205 від 21.02.2025 \par}
  {\large Формалізація вимог за стандартом ISO/IEC/IEEE 29148:2018 \par}
  \vfill
  {\large 2026 \copyright\ Системи Електронного Урядування \par}
\end{titlepage}

\newpage
\begin{center}
  {\Huge\color{primary}\bfseries Зміст\par}
\end{center}
\vspace{40pt}
\tableofcontents

\newpage
\section{Загальні відомості про засіб інформатизації}

\subsection{Найменування та підстави для розробки}
\textbf{Найменування засобу інформатизації:} Модуль «Істотність» (локальна ШІ-інфраструктура) у складі інформаційної системи управління підприємством ERP/1.

\textbf{Підстави для виконання робіт:}
\begin{itemize}
    \item Закон України «Про захист персональних даних»;
    \item Закон України «Про захист інформації в інформаційно-комунікаційних системах»;
    \item Порядок використання засобів інформатизації (Постанова КМУ № 205 від 21.02.2025);
    \item Програма модернізації обчислювальної IT-інфраструктури ЄСІКС.
\end{itemize}

\subsection{Призначення та цілі створення засобу}
Модуль призначений для розгортання локальних інтелектуальних сервісів аналізу та семантичного пошуку з використанням великих мовних моделей (LLM) та векторних баз даних (RAG) безпосередньо on-premises на базі споживчих GPU NVIDIA RTX 4070 Ti.

\textbf{Основні цілі створення:}
\begin{itemize}
    \item Забезпечення повного циклу локального інференсу LLM (Llama-3-8B-Instruct) із квантуванням ваг.
    \item Організація семантичного RAG-пошуку через RocksDB та FAISS.
    \item Автоматизація знеособлення вхідних документів (NER-маскування).
    \item Побудова оркестрованого контуру управління 600+ периферійними вузлами через K8s та Rancher за методологією GitOps.
    \item Реалізація відмовостійкості через CPU Fallback та Peer-to-Peer failover.
\end{itemize}

\subsection{План-графік виконання етапів робіт}
Розробка та впровадження модуля «Істотність» розділяється на такі етапи:
\begin{enumerate}
    \item \textbf{Етап 1: Технічне проектування} --- погодження UX/UI макетів, вибір цільових моделей ШІ, розробка схем даних. (1 місяць).
    \item \textbf{Етап 2: Реалізація сховища даних та RAG} --- налаштування RocksDB та індексів FAISS, розробка API-інтеграції з Erlang/OTP. (1.5 місяці).
    \item \textbf{Етап 3: Розробка NER та знеособлення} --- реалізація локального мікросервісу розпізнавання сутностей та маскування. (1 місяць).
    \item \textbf{Етап 4: Налаштування оркестрації K8s та GitOps} --- побудова кластера K8s, інтеграція GPU плагінів, оновлень. (1.5 місяці).
    \item \textbf{Етап 5: Державна експертиза КСЗІ} --- аудит безпеки, перевірка контролів NIST SP 800-53 та отримання атестату. (2 місяці).
\end{enumerate}

\newpage
\section{Відомості про робочий процес та умови експлуатації}

\subsection{Опис бізнес-процесів (BPMN / FSM)}
Всі процеси керування та роботи ШІ-сервісів формалізуються у вигляді кінцевих автоматів (FSM) та виконуються рушієм процесів BPE:

\subsubsection{Процес локального інференсу та RAG (ai\_inference.erl)}
Цей процес відповідає за обробку запиту користувача до великої мовної моделі з додаванням RAG-контексту:
\begin{enumerate}
    \item \textbf{StartInference:} Користувач надсилає текстовий запит (промпт) або документ.
    \item \textbf{AnonymizePrompt:} Фоновий процес NER-фільтрації знеособлює чутливі дані, замінюючи їх на токени.
    \item \textbf{GenerateEmbeddings:} Створення векторного представлення (ембедінгу) запиту через модель \texttt{e5-large}.
    \item \textbf{VectorSearch:} Пошук відповідних фрагментів документів у FAISS за косинусною близькістю.
    \item \textbf{BuildRAGPrompt:} Формування підсумкового промпту для LLM, де знеособлений запит поєднується зі знайденим контекстом.
    \item \textbf{ExecuteInference:} Надсилання промпту до локального сервера \texttt{llama.cpp} на GPU RTX 4070 Ti.
    \item \textbf{DeanonymizeResponse:} Зворотне відновлення персональних даних у згенерованій відповіді на базі локальної RocksDB таблиці відповідностей.
    \item \textbf{InferenceCompleted:} Завершення процесу, повернення відповіді користувачу.
\end{enumerate}

\subsubsection{Процес нічного оновлення моделей (ai\_deployment.erl)}
Забезпечує автоматизоване GitOps-оновлення ваг моделей:
\begin{enumerate}
    \item \textbf{TriggerUpdate:} Спрацювання cron-задачі о 01:00.
    \item \textbf{CheckRemoteCatalog:} Перевірка наявності нових версій моделей у центральному сховищі Scality RING.
    \item \textbf{VerifyDiskSpace:} Перевірка вільного місця на NVMe накопичувачі локального сервера.
    \item \textbf{PullModelWeights:} Асинхронне скачування GGUF файлу ваг із обмеженням швидкості.
    \item \textbf{VerifyChecksum:} Перевірка цілісності файлу за MD5/SHA256 хешами.
    \item \textbf{RolloutBlueGreen:} Оновлення поду в K8s зі зміною монтування на новий файл моделі.
    \item \textbf{VerificationTest:} Виконання тестового інференсу, перевірка GPU-метрик.
    \item \textbf{DeploymentCompleted:} Фіксація успішного оновлення.
\end{enumerate}

\subsection{Опис ролей та прав доступу (ABAC)}
Доступ до управління периферійними вузлами ШІ та перегляду журналів інференсу обмежується на базі ABAC:
\begin{verbatim}
allow(User, Action, Node) ->
    UserRole = maps:get(role, User),
    UserRegion = maps:get(region, User),
    NodeRegion = maps:get(region, Node),
    
    IsCDTO = (UserRole == cdto_admin),
    MatchesRegion = (UserRegion == NodeRegion),
    
    case {IsCDTO, Action} of
        {true, _} -> true;
        {false, view_metrics} -> MatchesRegion;
        {false, query_inference} -> true;
        _ -> false
    end.
\end{verbatim}

\subsection{Специфікація апаратного забезпечення}

В якості периферійного апаратного прискорювача використовуються споживчі GPU NVIDIA GeForce RTX 4070 Ti або RTX 4070 Ti Super. Детальні технічні параметри прискорювачів наведено в таблиці:

\begin{center}
\small
\begin{longtable}{p{4cm} p{5cm} p{5cm}}
\toprule
\textbf{Параметр} & \textbf{RTX 4070 Ti} & \textbf{RTX 4070 Ti Super} \\
\midrule
Об'єм VRAM & 12 GB GDDR6X & 16 GB GDDR6X \\
Шина пам'яті & 192-bit & 256-bit \\
Пропускна здатність & 504 GB/s & 672 GB/s \\
CUDA-ядра & 7680 & 8448 \\
Тензорні ядра & 240 (4-gen) & 264 (4-gen) \\
Енергоспоживання & 285 W & 285 W \\
\bottomrule
\end{longtable}
\end{center}

\subsection{Умови експлуатації та системне середовище}
Робочі станції ШІ-серверів експлуатуються в приміщеннях канцелярій чи комутаційних шафах установ. Завдяки Midi-Tower корпусу та активному повітряному охолодженню RTX 4070 Ti, додаткове кондиціонування не є обов'язковим. Системне середовище включає Ubuntu LTS 24.04, драйвери NVIDIA 550+, CUDA 12.x та Docker/K8s runtime.

\subsection{Специфікація мовних моделей для локальної генерації}

Для забезпечення локальної генерації текстів та семантичного аналізу документів в межах апаратних обмежень споживчих GPU (з обсягом відеопам'яті 12 або 16 ГБ VRAM) використовуються квантовані великі мовні моделі (LLM) класів Llama та Phi:

\begin{enumerate}
    \item \textbf{Llama-3-8B-Instruct} (квантування Q4\_K\_M або Q8\_0):
    \begin{itemize}
        \item Використовується для загальних задач текстової генерації, підготовки проектів рішень та семантичного аналізу.
        \item При квантуванні Q4\_K\_M обсяг пам'яті для моделі разом із контекстом 8K складає близько 6.5 ГБ VRAM, що дозволяє паралельно запускати кілька сесій інференсу на GPU з 12 ГБ VRAM.
        \item При квантуванні Q8\_0 обсяг пам'яті становить близько 9 ГБ VRAM (рекомендовано для прискорювачів з 16 ГБ VRAM).
    \end{itemize}
    \item \textbf{Phi-3-medium-14B} (квантування Q4\_K\_M):
    \begin{itemize}
        \item Використовується для складного логічного аналізу та перевірки фактів.
        \item Потребує близько 10.5 ГБ VRAM для інференсу з контекстом 4K. Рекомендовано для прискорювачів з 16 ГБ VRAM.
    \end{itemize}
    \item \textbf{Phi-3-mini-4B} (квантування Q8\_0):
    \begin{itemize}
        \item Надлегка і швидка модель для задач експрес-класифікації та локального розпізнавання сутностей (NER).
        \item Потребує близько 4.8 ГБ VRAM, забезпечуючи максимальну швидкість інференсу (понад 100 токенів/сек) на обох варіантах прискорювачів.
    \end{itemize}
\end{enumerate}

\subsection{Архітектура локальної генерації та оптимізації VRAM}

Локальний інференс та генерація текстів на периферійних вузлах реалізуються за допомогою спеціалізованого оптимізованого інференс-сервера \texttt{llama.cpp}, запущеного в ізольованому Docker-контейнері під управлінням K8s. Взаємодія з графічним прискорювачем здійснюється через драйвер NVIDIA CUDA 12.x з використанням обчислювальних ядер CUDA та Tensor Cores 4-го покоління.

\subsubsection{Профіль розподілу та бюджетування VRAM (Memory Math)}

Для забезпечення безперебійної роботи без виходу за межі пам'яті (Out of Memory - OOM) на споживчих картах розраховано детальний бюджет відеопам'яті:

\begin{enumerate}
    \item \textbf{Сценарій 1: Робоча станція з NVIDIA RTX 4070 Ti (12 GB VRAM)}
    \begin{itemize}
        \item \textit{Цільова модель}: Llama-3-8B-Instruct (квантування Q4\_K\_M, вага файлу $\approx$ 4.8 GB).
        \item \textit{Статичне резервування під модель}: $\approx$ 5.2 GB VRAM (з урахуванням завантаження CUDA графів).
        \item \textit{KV-кеш (Context Window 8K)}: Об'єм пам'яті під KV-кеш для одного активного слота користувача при $N_{layers}=32$, $N_{heads}=8$, $D_{head}=128$ та довжині контексту $8192$ токенів розраховується як:
        $$\text{Size}_{KV} = 2 \times N_{layers} \times N_{heads} \times D_{head} \times 2 \text{ bytes (FP16)} \times 8192 = 1.07 \text{ GB}$$
        \item \textit{Багатокористувацький ліміт}: Для підтримки 5 паралельних користувацьких сесій виділяється $5 \times 1.07\text{ GB} = 5.35\text{ GB}$ VRAM.
        \item \textit{Загальний обсяг використання VRAM}: $5.2\text{ GB} (\text{модель}) + 5.35\text{ GB} (5\text{ сесій KV}) + 1.0\text{ GB} (\text{системний оверхед/драйвери}) = 11.55\text{ GB}$ VRAM (запас $\approx$ 450 MB).
    \end{itemize}

    \item \textbf{Сценарій 2: Робоча станція з NVIDIA RTX 4070 Ti Super (16 GB VRAM)}
    \begin{itemize}
        \item \textit{Цільова модель}: Phi-3-medium-14B (квантування Q4\_K\_M, вага файлу $\approx$ 8.5 GB).
        \item \textit{Статичне резервування під модель}: $\approx$ 9.1 GB VRAM.
        \item \textit{KV-кеш (Context Window 4K)}: Об'єм пам'яті під KV-кеш для одного користувача при контексті $4092$ токенів становить $\approx$ 0.75 GB VRAM.
        \item \textit{Багатокористувацький ліміт}: Для підтримки 8 паралельних сесій виділяється $8 \times 0.75\text{ GB} = 6.0\text{ GB}$ VRAM.
        \item \textit{Загальний обсяг використання VRAM}: $9.1\text{ GB} (\text{модель}) + 6.0\text{ GB} (8\text{ сесій KV}) + 0.8\text{ GB} (\text{оверхед CUDA}) = 15.9\text{ GB}$ VRAM.
    \end{itemize}
\end{enumerate}

\subsubsection{Паралельне планування та оптимізація пропускної здатності}

Для досягнення високої швидкості генерації та ефективного використання Tensor Cores застосовуються такі технологічні механізми:
\begin{itemize}
    \item \textbf{Continuous Batching (Ітераційне пакетування)}: Інференс-сервер \texttt{llama.cpp} запускається з параметром \texttt{--cont-batching} (або \texttt{-cb}), що дозволяє обробляти запити в режимі реального часу. Нові запити додаються в конвеєр обробки на кожній ітерації декодування токенів без необхідності очікування завершення попередньої генерації.
    \item \textbf{PagedAttention (Сторінковий KV-кеш)}: Пам'ять KV-кешу розбивається на логічні сторінки (блоки токенів розміром 16 чи 32) та динамічно виділяється через таблицю сторінок, аналогічно віртуальній пам'яті в ОС. Це повністю запобігає фрагментації відеопам'яті та дозволяє оптимізувати простір під паралельні запити.
    \item \textbf{FlashAttention / SDPA}: Використання оптимізованих ядер обчислення уваги, що значно знижує кількість операцій читання/запису у глобальну пам'ять GPU та прискорює етап Prefill у 2.5--3 рази.
\end{itemize}

\newpage

\section{Інформаційне та програмне забезпечення}

\subsection{Інформаційне забезпечення (Схеми та Дані)}

\subsubsection{Визначення структур даних (Erlang Records)}
Всі ключові сутності модуля описуються у вигляді рекорді Erlang:
\begin{verbatim}
-record(ai_node_status, {
    node_id,                    % UUID вузла (PK)
    hostname,                   % Текстове ім'я хоста
    ip_address,                 % IP-адреса вузла
    gpu_temp = 0,               % Температура GPU RTX 4070 Ti в градусах
    vram_used = 0,              % Зайнято відеопам'яті (MB)
    vram_free = 0,              % Вільно відеопам'яті (MB)
    power_limit_pct = 80,       % Power Limit у відсотках (80% за замовчуванням)
    status = offline,           % online | offline | degraded | cooling_down
    last_seen_at = 0            % Timestamp останнього пінгу
}).

-record(ai_model_config, {
    model_id,                   % UUID моделі (PK)
    name,                       % Назва моделі (напр. Llama-3-8B-Instruct)
    quant_type,                 % Тип квантування (q4_k_m | q8_0)
    file_path,                  % Шлях до GGUF файлу на NVMe
    checksum_sha256,            % Контрольна сума файлу
    context_window = 8192,      % Максимальний контекст (токенів)
    is_active = false,          % Чи завантажена модель у пам'ять GPU
    updated_at = 0
}).

-record(ai_anonymization_dictionary, {
    dict_id,                    % UUID словника (PK)
    mapping_table = #{},        % Таблиця відповідностей {Token -> OriginalValue}
    created_at = 0,
    updated_at = 0
}).

-record(ai_inference_task, {
    task_id,                    % UUID задачі (PK)
    user_id,                    % UUID користувача, що надіслав запит
    node_id,                    % UUID обчислювального вузла, який виконав запит
    prompt_hash,                % SHA256 хеш промпту
    response_hash,              % SHA256 хеш згенерованої відповіді
    tokens_generated = 0,       % Кількість згенерованих токенів
    time_taken_ms = 0,          % Час виконання запиту в мілісекундах
    status = pending,           % pending | processing | completed | failed
    created_at = 0
}).

-record(ai_rag_index_state, {
    index_id,                   % UUID індексу (PK)
    num_vectors = 0,            % Кількість векторів
    file_path,                  % Шлях до FAISS HNSW індексу
    last_sync_at = 0,           % Час останньої синхронізації з Scality RING
    status = active             % active | syncing | stale
}).
\end{verbatim}

\subsection{Програмне забезпечення (Реалізація)}

\subsubsection{ai\_inference.erl (DSL процесу RAG інференсу)}
\begin{verbatim}
-module(ai_inference).
-include("bpe.hrl").
-compile(export_all).

def() ->
    #process{
        name = 'AI Inference and RAG Workflow',
        beginEvent = 'StartInference',
        endEvent = 'InferenceCompleted',
        tasks = [
            #beginEvent{name='StartInference'},
            #serviceTask{name='AnonymizePrompt', module=?MODULE},
            #serviceTask{name='GenerateEmbeddings', module=?MODULE},
            #serviceTask{name='VectorSearch', module=?MODULE},
            #serviceTask{name='BuildRAGPrompt', module=?MODULE},
            #serviceTask{name='ExecuteInference', module=?MODULE},
            #serviceTask{name='DeanonymizeResponse', module=?MODULE},
            #endEvent{name='InferenceCompleted'}
        ],
        flows = [
            #sequenceFlow{name='1', source='StartInference', target='AnonymizePrompt'},
            #sequenceFlow{name='2', source='AnonymizePrompt', target='GenerateEmbeddings'},
            #sequenceFlow{name='3', source='GenerateEmbeddings', target='VectorSearch'},
            #sequenceFlow{name='4', source='VectorSearch', target='BuildRAGPrompt'},
            #sequenceFlow{name='5', source='BuildRAGPrompt', target='ExecuteInference'},
            #sequenceFlow{name='6', source='ExecuteInference', target='DeanonymizeResponse'},
            #sequenceFlow{name='7', source='DeanonymizeResponse', target='InferenceCompleted'}
        ],
        roles = [operator, system]
    }.

action({serviceTask, 'AnonymizePrompt'}, Proc) ->
    % Виклик локального мікросервісу знеособлення NER
    {reply, Proc};
action({serviceTask, 'VectorSearch'}, Proc) ->
    % Пошук у локальному індексі FAISS
    {reply, Proc};
action({serviceTask, 'ExecuteInference'}, Proc) ->
    % Інференс на GPU RTX 4070 Ti через llama.cpp
    {reply, Proc};
action(_, Proc) -> {reply, Proc}.
\end{verbatim}

\subsubsection{ai\_validator.erl (Модуль валідації ШІ вузла)}
\begin{verbatim}
-module(ai_validator).
-export([validate_model_hash/2, check_gpu_temperature/1]).

validate_model_hash(FilePath, ExpectedHash) ->
    case file:read_file(FilePath) of
        {ok, Binary} ->
            Calculated = crypto:hash(sha256, Binary),
            case Calculated of
                ExpectedHash -> ok;
                _ -> {error, checksum_mismatch}
            end;
        _ -> {error, file_not_found}
    end.

check_gpu_temperature(GpuTemp) ->
    MaxAllowed = 80,
    if
        GpuTemp > MaxAllowed -> {error, overheating};
        true -> ok
    end.
\end{verbatim}

\newpage
\section{Вимоги до засобу за ISO/IEC/IEEE 29148}

\subsection{Функціональні вимоги}
Модуль «Істотність» повинен виконувати аналіз завантажених документів, генерувати семантичні ембедінги, проводити знеособлення даних, здійснювати пошук контексту у RocksDB/FAISS та забезпечувати інференс мовної моделі.

\subsection{Вимоги до інтерфейсів та дизайну}
\begin{itemize}
    \item Веб-інтерфейс адміністратора вузла повинен базуватися на фреймворку N2O/NITRO.
    \item Панель моніторингу повинна виводити статус GPU (температура, VRAM) та кількість активних сесій.
    \item Дизайн інтерфейсів має повністю відповідати стилю ERP/1 (Sleek Dark Mode).
\end{itemize}

\subsection{Вимоги до безпеки та захисту інформації}
\begin{itemize}
    \item Шифрування розділів NVMe накопичувачів на базі LUKS з інтеграцією апаратного модуля TPM 2.0.
    \item Ізоляція мережевого трафіку ШІ від загального контуру LAN за допомогою Network Policies в CNI-плагіні Cilium.
    \item Заборона передачі будь-яких даних за межі локальної фізичної мережі установи (Air-Gapped режим).
\end{itemize}

\subsection{Вимоги до продуктивності та надійності}
\begin{itemize}
    \item Середній показник доступності (SLA) периферійного сервера не менше 99.9\%.
    \item Час відповіді на RAG-пошук (векторизація + FAISS) не більше 250 мс.
    \item Максимальний час перемикання на CPU при виході з ладу GPU не більше 15 секунд.
\end{itemize}

\subsection{Вимоги до логування та аудиту}
\begin{itemize}
    \item Логування кожного запиту: ідентифікатор користувача, час, хеш промпту, хеш відповіді, згенеровані токени.
    \item Журнал аудиту повинен бути захищений від змін та автоматично передаватися на центральний сервер логування в нічний час.
\end{itemize}

\subsection{Адміністративні та юридичні вимоги}
\begin{itemize}
    \item Всі компоненти ПЗ поставляються з відкритими ліцензіями (MIT/Apache 2.0).
    \item Застосування споживчих GPU RTX 4070 Ti у судах та офісних приміщеннях не підпадає під заборону використання GeForce у ЦОД (Datacenters) згідно з EULA NVIDIA, оскільки вузли є робочими станціями (Workstations) або крайовими пристроями (Edge Devices).
\end{itemize}

\end{document}
